% --- Archon physics formalization source begin ---
% archon:physics
% archon:covers IPhO2026Problems/problem_ipho_2026_e1_c7.lean
% archon:source-report reports/ipho_2026/problem_ipho_2026_e1_c7.source.json
% archon:problem-id ipho_2026_e1
% archon:part-id E1-C7
% archon:previous-part ipho_2026_e1_c6
% archon:previous-part-policy derive_inline_from_problem_only_material

\chapter{Ipho Physics problem ipho\_2026\_e1\_c7}
\label{ch:IPhO2026Problems_problem_ipho_2026_e1_c7}

\paragraph{Problem source.}
\#\# Physical scenario

Water in the inner and outer cylinders exchanges heat radially through an
acrylic cylindrical wall.  Record T\_IC and T\_OC versus time.  The heat-flow
model is dQ/dt = (T\_OC - T\_IC)/R\_Th.  For radial Fourier conduction,
dQ/dt = -lambda*A*dT/dr.  Ignore apparatus heat capacity where instructed and
use the dimensions in Figure 17.

\#\# Current subquestion E1-C7

Combine the heat-flow relation and radial Fourier law to determine acrylic conductivity lambda.

\paragraph{Shared problem context.}
Water in the inner and outer cylinders exchanges heat radially through an
acrylic cylindrical wall.  Record T\_IC and T\_OC versus time.  The heat-flow
model is dQ/dt = (T\_OC - T\_IC)/R\_Th.  For radial Fourier conduction,
dQ/dt = -lambda*A*dT/dr.  Ignore apparatus heat capacity where instructed and
use the dimensions in Figure 17.

\paragraph{Current subquestion.}
Combine the heat-flow relation and radial Fourier law to determine acrylic conductivity lambda.

\paragraph{Figure/image paths.}
\begin{itemize}
\item ipho\_2026\_source/image/E1\_page-14.png
\item ipho\_2026\_source/image/E1\_page-13.png
\end{itemize}

\paragraph{Reusable previous-part conclusions.}
\begin{itemize}
\item Source E1-C6. Question: Determine the effective wall thermal resistance R\_Th from the C5 graph. Policy: derive\_inline\_from\_problem\_only\_material
\end{itemize}

\paragraph{Formalization target.}
create a compiling Lean file with sorry bodies at `IPhO2026Problems/problem\_ipho\_2026\_e1\_c7.lean`.
The Lean declarations must preserve the physical quantities, dimensions or dimensional roles, figure labels, governing-law hypotheses, and final relation expressed by this problem.
Use Mathlib/Physlib names found through LeanExplore where available. If a domain API is missing, introduce faithful local abstractions rather than scalar placeholder aliases.
This is an answer-blind solve-phase task. Derive a candidate result only from the problem-side evidence above; do not assume a target value, select a tolerance around a desired value, or encode a desired result into a candidate domain.
For numerical questions, define the raw end-to-end quantity and apply an explicit source-derived rounding rule. For identification questions, characterize or prove uniqueness before recording the derived candidate.

\begin{theorem}[Ipho Physics formalization target]
\label{thm:physics:ipho_2026_e1_c7:target}
The assigned autoformalize agent should translate this IPhO physics problem into Lean declarations in the covered file, with theorem and lemma proof bodies written as `by sorry`.
\end{theorem}
\begin{proof}
This is an autoformalization task, not a proof task. Produce faithful statements that can later be proved without weakening the source contract.
\end{proof}
% --- Archon physics formalization source end ---
